% !TeX root = ../../Book3.tex
\chapter{共形映射的边界行为}
\label{b3:ch:07}
\BookChapterNumber{b3:ch:07}
\begin{chapterabstract}
本章从内部共形映射走向边界。横割与素端记录不同的边界趋近方式，Jordan 区域的边界同胚为多边形公式提供预顶点；直边反射与角度幂次导出 Schwarz--Christoffel 公式。矩形的边长计算引出完全椭圆积分。一般素端紧化和局部连通判别作为明确标注的深层输入，其应用证明与具体计算完整展开。
\end{chapterabstract}

\begin{learninggoals}
\begin{itemize}
\item\label{b3:item:07-goal-1}区分普通边界点、素端和素端印象，识别裂缝两侧的不同通道。
\item\label{b3:item:07-goal-2}准确使用连续延拓与边界同胚各自的假设。
\item\label{b3:item:07-goal-3}从反射和内角推导 Schwarz--Christoffel 导数公式。
\item\label{b3:item:07-goal-4}用椭圆积分确定矩形的形状参数，并核对其方向约定。
\end{itemize}
\end{learninggoals}

第5章的圆盘映射把区域内部的点一一对应，却没有给边界上的坐标。对于裂平面，同一个裂缝点可由两侧到达；这两侧在圆盘上应占不同位置。边界行为的任务因此先是寻找合适的边界对象，然后才是计算映射。

\ProseParagraphBreak
本章所用的内部分析前提是 Riemann 映射、Morera 定理、辐角原理与局部标准形。一般素端对应涉及较深平面拓扑，证明状态在定理旁明确说明；在此输入上，多边形与矩形的计算不再隐藏额外边界假设。

% !TeX root = ../../Book3.tex
\section{横割与素端}
\label{b3:sec:0701}

内部共形坐标不自动给出边界坐标。裂缝两侧在平面中是同一组点，却对应不同的趋近方式。横割与素端用区域内部的分离来记录这种方式。本节先在有界单连通区域 \(U\) 中定义；无界情形在球面中取边界、闭包与直径，允许无穷远作为端点；这给出相应的不依赖有限坐标的版本。

\subsection{横割}
\label{b3:subsec:070101}

\begin{definition}\label{b3:def:07-crosscut}
区域 \(U\) 的\term[hengge]{横割}[crosscut]是一个开 Jordan 弧 \(C\subset U\)，它的闭弧有两个不同的端点在 \(\partial U\)，其余点都在 \(U\)。
\end{definition}
对于单连通平面区域，横割把 \(U\) 分成恰好两个连通分支。这是 Jordan 分离定理在横割上的形式：沿区域内路径把弧的两侧闭合，得到一个 Jordan 圈；两侧不能绕过边界端点相通。圆盘中的弦是基本例子，弧可以弯曲而不必是直线。若弧的两个端点在内部，它通常不分离区域；若闭弧自交，也不能套用横割分离结论。

\ProseParagraphBreak
固定观察点 \(a\in U\)。当 \(a\notin C\) 时，记不含 \(a\) 的横割分支为 \(D(C)\)。它代表从 \(a\) 看被横割挡住的一侧。边界接近的定义将使用这些内部开集，而不是先假定每个边界点都有一条可达路径。

\ProseParagraphBreak
横割的两个端点属于边界，内部属于区域；这三项条件都参与分离结论。单位盘中的弦是最简单的模型：直弦把圆盘分成两个半盘型区域，圆周上两个开弧分别贴近两个分支。若把弦改成只以一个端点落在边界的线段，它的内部终点可以绕过去，不再把圆盘分开。若把两端都落在区域内部，也没有相同的两侧分离结论。一般单连通域的横割分离是平面 Jordan 拓扑的应用，并不是单凭曲线连续就能推出。

\ProseParagraphBreak
选定观察点的作用是给横割的两侧作一致选择。若观察点落在横割上，先把观察点换到区域内另一个点；链趋近边界后，两点最终处在同一个远侧分支，因此不会改变链的端部等价类。这里的“最终”很关键：前面有限几条横割的选边可以不同，但素端由链的尾部决定。把有限项删除或重新编号，不应制造新的素端。

\subsection{链}
\label{b3:subsec:070102}

\begin{definition}\label{b3:def:07-null-chain}
一列横割 \(C_n\) 称为一条零链，若闭横割两两不交，\(\operatorname{diam}C_n\rightarrow0\)，且在舍去有限项后，\(D_n=D(C_n)\) 满足 \(D_{n+1}\subset D_n\)、\(C_{n+1}\subset D_n\)，同时 \(\bigcap_nD_n=\varnothing\)。
\end{definition}
小直径约束的是横割而非整个 \(D_n\)。复杂边界中一个很短的横割后面仍可能藏着一个很长的狭窄通道。链的印象定义为
\[
 I=\bigcap_n\overline{D_n}.
\]
在有界情形，它是非空紧连通集且包含于边界。非空来自嵌套紧集，连通性来自连通闭包的嵌套交；若交内有内点，则该点的小圆盘与充分小、远离此点的横割不交，连通性迫使它一直处在同一侧，最终与内部交为空矛盾。不能仅凭横割直径趋于0就把印象写成单点。

\ProseParagraphBreak
对于圆盘，固定边界点附近的短弦给出一条链，分支为越来越小的圆帽，印象恰为该点。对于裂盘，同一个裂缝内点有从上方与下方趋近的两条链；它们的印象可能相同，却不代表同一个边界通道。

\subsection{素端}
\label{b3:subsec:070103}

若对每个 \(n\)，另一条链的充分晚分支都包含在 \(D_n\) 中，称后一链从属前一链。两链互相从属时称为等价。传递性由“先给第一链的一个分支，再依次选择后两链的充分晚分支”得到，所以等价是等价关系。

\begin{definition}\label{b3:def:07-prime-end}
零链的等价类称为 \(U\) 的\term[suduan]{素端}[prime end]。其印象是任一代表链的闭分支交。
\end{definition}
互相从属保证两个代表给出的印象相同。改变观察点也不影响充分晚的从属关系：连接两个观察点的固定内部紧路径不会与充分晚的横割相交，两个点最终落在同一侧。定义依靠尾部，故舍掉有限条横割不改变素端。

\ProseParagraphBreak
“素端”与“边界点”属于不同集合。印象是一个集合，素端是趋近通道的等价类；既可能一个素端有非单点印象，也可能不同素端有相同单点印象。对有规则边界的区域，二者恰巧一致；这需要边界对应定理，不能写入一般定义。

\subsection{素端边界}
\label{b3:subsec:070104}

令 \(\widehat U^{\mathrm P}=U\cup\partial_{\mathrm P}U\)。内部沿用通常拓扑；由一个链分支 \(D_n\) 加上所有从属于该分支的素端，定义该素端的一组邻域。互相从属使这种拓扑与代表链无关。这一空间称为素端紧化，其边界是 \(\partial_{\mathrm P}U\)。

\begin{theorem}[深层素端对应定理]\label{b3:thm:07-prime-end-correspondence}
设 \(U\subsetneq\mathbb C\) 单连通，\(f:\mathbb D\rightarrow U\) 双全纯。则 \(f\) 唯一延为 \(\overline{\mathbb D}\) 到 \(\widehat U^{\mathrm P}\) 的同胚。特别地，素端边界具有一个圆周的拓扑与循环次序。
\end{theorem}
\begin{proofsketch}
这个定理包含一般平面边界的分离与链逼近，不在这里声称完成全部拓扑证明。路线是用横割分支的嵌套关系建立素端拓扑；再用共形映射的面积估计，选取趋向圆周的短弧，使其像形成零链；反向以圆盘内的横割分离证明同一素端只能对应一个圆周点；最后核对分支邻域对应圆帽邻域。困难在于普通边界可能并非局部连通，不能假定横割像的闭弧已有边界端点。完整处理可读 Lars V. Ahlfors 的《Conformal Invariants》\cite{Ahlfors2010Conformal}的 \S4-6（57--60页），其中也说明横割链与圆盘边界的对应；Christian Pommerenke 的《Boundary Behaviour of Conformal Maps》\cite{Pommerenke1992Boundary}第2章则为进一步阅读。这里保留该定理为明确的深层输入。
\end{proofsketch}
这个输入给边界一个合适的拓扑对象，并没有保证 \(f\) 在每个圆周点有唯一普通边界极限。由素端印象到平面边界的投影可能有重合或非单值集合像。下节说明何时能把这项结果化为通常的连续延拓。

\ProseParagraphBreak
裂缝圆盘提供必须保留两侧的例子。把单位盘沿正实半径切开，在裂缝内部的一个边界点 \(x\in(0,1)\)，可以从上侧和下侧分别取越来越小的半圆横割。两条链的印象都是 \(\{x\}\)，但选出的近侧区域不最终互相包含：一条始终位于裂缝上侧，另一条始终位于下侧。因此它们对应两个不同素端。素端记录趋近方向，印象只记录最后落在哪些平面边界点，二者不能一一混同。裂缝端点附近的链又可能把两侧接在同一个角端，须看局部横割的实际分离。

\ProseParagraphBreak
在圆盘中，边界点 \(\zeta\) 的标准链可取 \(\mathbb D\cap\{|z-\zeta|=\varepsilon_n\}\)，其中 \(\varepsilon_n\downarrow0\)。选靠近 \(\zeta\) 的分支，其闭包交为 \(\{\zeta\}\)。链邻域与通常的圆盘闭包邻域相容。映射到一般单连通域后，平面直径未必仍按同样速度趋于0，素端对应定理正是把这种内部对应和边界拓扑协调起来；它的证明不能仅用连续映射将小圆变成小曲线的一句话代替，因为此时边界连续性尚未建立。

% !TeX root = ../../Book3.tex
\section{Carathéodory 边界对应}
\label{b3:sec:0702}

Jordan 区域的边界是一条简单闭曲线，此时每个边界点只对应一个素端。若边界仅局部连通，则可得到连续延拓，但边界映射仍可能多对一。三种陈述应当分别理解。

\subsection{Jordan 区域}
\label{b3:subsec:070201}

\term[jordan-quyu]{Jordan 区域}[Jordan domain]指以 Jordan 曲线为边界的有界内部区域。Jordan 曲线是一张圆周的连续单射像，不要求可微、有限长度或有切线。曲线的紧性使该参数化自动是同胚。多边形内部、圆盘及具有尖点但没有自接触的简单闭曲线内部都属于此类。

\ProseParagraphBreak
平面分离定理进一步给 Jordan 曲线两侧的拓扑圆盘结构。在 Jordan 区域中，以靠近一个边界点的两段边界弧和内部短弧构成横割，便能得到印象只含该点的链；反过来，每条零链的分支在边界圆周次序中挡住越来越短的弧。因此素端与普通边界点一一对应。这个平面拓扑事实不涉及全纯函数的导数，也不要求边界整齐。

\subsection{连续边界延拓}
\label{b3:subsec:070202}

\begin{theorem}[深层连续延拓判别]\label{b3:thm:07-continuous-extension}
设 \(U\) 有界且单连通，\(f:\mathbb D\rightarrow U\) 双全纯。则 \(f\) 有到 \(\overline{\mathbb D}\) 的连续延拓，当且仅当 \(\partial U\) 局部连通。
\end{theorem}
\begin{proofsketch}
本定理作为素端理论的进一步输入。局部连通性使每个素端印象为单点，并使印象投影连续，于是与素端对应复合得到延拓。反向利用连续圆周像的局部连通性；圆周参数化的连续性先给均匀小像，再按有限弧分割控制邻域中的连通片。一般边界上的精确拓扑论证见前引 Pommerenke 第2章。此处不把仅有印象单点的点态结论误当成全圆周连续性证明。
\end{proofsketch}
裂盘的边界局部连通，因此映射可连续延到圆周，但裂缝两侧对应不同的圆周点。连续延拓可以满射到整个边界，却不必单射。无界区域应在球面中作相应陈述；欧氏意义的有限边界值可能在无穷远点失效。

\subsection{边界同胚}
\label{b3:subsec:070203}

\term[caratheodory-bianjie-dingli]{Carathéodory 边界对应定理}[Carathéodory boundary correspondence theorem]在 Jordan 情形给出一一对应。

\begin{theorem}[Carathéodory 边界同胚]\label{b3:thm:07-caratheodory-jordan}
若 \(U\) 为 Jordan 区域，任何双全纯 \(f:\mathbb D\rightarrow U\) 唯一延为 \(\overline{\mathbb D}\rightarrow\overline U\) 的同胚。
\end{theorem}
\begin{proof}
Jordan 边界的素端投影是同胚，如本节第一小节的平面分离描述。将它与定理\ref{b3:thm:07-prime-end-correspondence}的素端延拓复合，得到闭圆盘到闭区域的连续双射。源空间紧、目标 Hausdorff，因此逆也连续。内部延拓的唯一性由内部点稠密得到。
\end{proof}
这是一条依赖已声明深层输入的完整推论证明，而不是对全部素端理论的替代。它不声称边界上的导数存在。例如角点附近的导数一般按一个幂次趋于0或无穷大；边界同胚仍然完全成立。

\begin{proposition}[圆边反射的局部形式]\label{b3:prop:07-circle-reflection}
若全纯函数 \(g\) 在圆的一侧邻域内无零，连续到一段圆弧，且该弧上 \(|g|=1\)，则可按反射公式跨弧延拓。对于单位圆，另一侧的值为
\[
 \widetilde g(z)=\frac1{\overline{g(1/\bar z)}}.
\]
\end{proposition}
\begin{proof}
反演 \(z\mapsto1/\bar z\) 交换单位圆两侧，复共轭和倒数把原侧全纯无零函数变成另一侧的全纯函数；圆弧上 \(|g|=1\) 使两个表达式相等。局部用一个 Möbius 变换把该圆弧拉直，再对目标单位圆作同样的拉直，转化为直径上取实值的反射。直边反射给拼接全纯，逆变换恢复所写公式。选择的小邻域避开 Möbius 变换的极点，且无零条件保证倒数合法。
\end{proof}
若圆弧是边界的一部分，这只延拓到该弧的局部另一侧，不保证映射仍在整个扩大区域单射。\(g(z)=z^m\) 的单位圆边界值模为1，确可全纯反射，但 \(m>1\) 时在盘内不是单叶映射。反射延拓与全局单射是两种不同结论。

\ProseParagraphBreak
光滑弧上的反射还给边界导数的意义。若源和目标边界在一点附近均为解析弧，且映射作为边界同胚把两段对应，拉直后得到两侧全纯映射。只要延拓后的零点阶数为1，就有非零边界导数；若目标有角点，则需先用幂次拉直，导数可趋于0或无穷。因此“边界延拓连续”不足以推出边界导数存在，更不足以推出导数不为零。

\subsection{素端表述}
\label{b3:subsec:070204}

素端表述在作边界标记时尤其有用。一个四边形需指定四个循环有序的边界素端，而不能仅给四个平面点；在裂缝区域，同一个点可有两个不同的标记。边界对应把这种标记传到圆盘圆周，随后可由 Möbius 变换把三个标记固定。剩下的第四个位置给出一个实际的共形参数。

\ProseParagraphBreak
若只关心内部映射，第5章不需要任何素端前提；若需要多边形边界值或矩形四角位置，就必须使用本章的边界同胚。这个依赖界面解释了为何 Schwarz--Christoffel 公式的证明虽然从一个内部 Riemann 映射开始，却还需要辨认实轴上的预顶点。

\ProseParagraphBreak
边界映射的三个层次可以在计算前作一次检查。第一层只有素端同胚，得到圆周参数，但一个参数点可能对应普通边界中的一整个印象。第二层有普通闭包上的连续映射，参数点有唯一边界值，但不同参数点仍可能对应同一点。第三层是 Jordan 边界同胚，参数点与普通边界点一一对应。Schwarz--Christoffel 的预顶点需要第三层，裂缝边界的两侧计算至少要保留第二层中的两份参数资料。若把裂缝的两岸在预顶点列表里合成一条边，就会丢掉循环次序。

\ProseParagraphBreak
Jordan 边界没有隐含弧长参数。某条简单闭曲线可以长度无穷，Carathéodory 定理仍给同胚，但沿它的分段光滑积分未必有定义。要写经典边界 Cauchy 积分，还需边界可求长或改用内部逼近围道。反射则需要解析边段，角点公式需要确切的局部角域；这些要求都比边界同胚更强。具体应用应当先决定所需的是拓扑值、路径积分还是导数，再列对应的边界条件。

% !TeX root = ../../Book3.tex
\section{Schwarz–Christoffel 公式}
\label{b3:sec:0703}

多边形的边界曲率集中在顶点。直边允许 Schwarz 反射，顶点内角决定映射导数的幂次。把这些局部信息合并，就得到可积分的全局公式。

\subsection{多边形区域}
\label{b3:subsec:070301}

设 \(P\) 是有界简单多边形内部，顶点按正向顺序为 \(w_1,\ldots,w_n\)，内角为 \(\alpha_j\pi\)、\(0<\alpha_j<2\)。由 Riemann 映射与 Jordan 边界同胚，可把上半平面双全纯映到 \(P\)，并在扩充实轴上标记预顶点 \(x_j\)。选取上半平面自同构，可令无穷远对应某条边的内部点，故所有 \(x_j\) 都有限。

\term[schwarz-fanshe-yuanli]{Schwarz 反射原理}[Schwarz reflection principle]先在直边上表述。

\begin{proposition}[直边反射]\label{b3:prop:07-schwarz-reflection}
若 \(g\) 在上半圆盘全纯，连续到直径且直径上取实值，则它按 \(g(\overline z)=\overline{g(z)}\) 全纯延到整个圆盘。
\end{proposition}
\begin{proof}
定义下半盘上的反射式，直径实值使拼接连续。任何跨直径的小三角形可沿距直径 \(\varepsilon\) 的平行线分割，两侧积分为零；令 \(\varepsilon\downarrow0\)，新增切边相消。Morera 定理给拼接函数全纯。对非水平直边先作仿射旋转即可。
\end{proof}
反射要求直径上确实有连续实值。只知道上半盘全纯、或仅知道几乎处处实边界值，不足以直接套用这项初等版本。

\subsection{内角}
\label{b3:subsec:070302}

在 \(w_j\) 附近用 \((w-w_j)^{1/\alpha_j}\) 把角域拉直到半平面。分支选在多边形局部内部，连续到两条边。与 \(f\) 复合后，可跨 \(x_j\) 反射。拉直后的映射在 \(x_j\) 有一个简单零点：若阶数大于1，局部标准形将在半邻域中覆盖过大的角度，违反其到半平面的单射性与两侧对应。故有
\[
 f(z)-w_j=(z-x_j)^{\alpha_j}h_j(z),\qquad h_j(x_j)\ne0,
\]
其中幂函数在上半邻域取分支，\(h_j\) 在小邻域全纯无零。于是
\[
 \frac{f''(z)}{f'(z)}=\frac{\alpha_j-1}{z-x_j}+\text{全纯项}.
\]
凸角 \(\alpha_j<1\) 使导数发散，凹角 \(\alpha_j>1\) 使导数趋于0。这都不妨碍内部导数非零，因为顶点在边界上。

\ProseParagraphBreak
角点指数可以直接在扇形中核验。令 \(W_\alpha=\{0<\arg w<\alpha\pi\}\)，在上半平面取 \(w=z^\alpha\)。两条实半轴分别映到扇形两边；正实侧的导数为 \(\alpha x^{\alpha-1}\)。当 \(\alpha=1/2\)，直角的局部映射是平方根，导数发散如 \(x^{-1/2}\)，但函数本身连续且差值为 \(O(x^{1/2})\)。当 \(\alpha=3/2\)，凹角的局部差值为 \(O(x^{3/2})\)，导数趋于0。把角点拉直后，导数又恢复有限非零。这说明坐标必须随几何选择，不能把某个坐标中的导数大小作为不变的边界奇性。

\ProseParagraphBreak
若角域只在有限邻域内成立，幂次必须在那一邻域的内部选定。不能在环绕顶点的整个穿孔平面上不加分支约定地写 \((w-w_j)^{1/\alpha_j}\)。在实际多边形中，两条邻边夹出的内部角恰提供所需的连续辐角范围，因而局部幂次合法。余下无零因子可在小盘取对数，其 \(\alpha_j\) 次幂吸收到 \(h_j\) 中，得到正文的角点标准形。

\subsection{公式}
\label{b3:subsec:070303}

\term[schwarz-christoffel-gongshi]{Schwarz--Christoffel 公式}[Schwarz--Christoffel formula]把上述角点资料合并成映射导数。

\begin{theorem}[Schwarz--Christoffel 公式]\label{b3:thm:07-schwarz-christoffel}
在上述所有预顶点有限的归一化中，存在 \(A\in\mathbb C\)、\(C\ne0\)，使
\[
 f(z)=A+C\int_{z_0}^{z}\prod_{j=1}^n(\zeta-x_j)^{\alpha_j-1}\,\mathrm{d}\zeta.
\]
各幂在上半平面连续取分支，积分与内部路径无关。
\end{theorem}
\begin{proof}
在每个非顶点实区间，\(f'\) 的辐角是所对应边的固定方向，所以 \(f''/f'\) 取实值，并可反射到下半平面。顶点处前一小节给简单极点及留数 \(\alpha_j-1\)。无穷远对应边内点，使用坐标 \(t=1/z\) 的反射与局部单射性得 \(f'(z)=c z^{-2}(1+O(1/z))\)，故 \(f''/f'=-2/z+O(1/z^2)\)。多边形角和给 \(\sum_j(\alpha_j-1)=-2\)。减去 \(\sum_j(\alpha_j-1)/(z-x_j)\) 后，差是球面上无极点、在无穷远趋于0的函数，故由 Liouville 为0。积分对数导数，得到 \(f'=C\prod_j(z-x_j)^{\alpha_j-1}\)。上半平面单连通且乘积全纯，再取原函数得公式。
\end{proof}
这项证明从已存在的多边形映射推出公式。任意指定角度和预顶点的积分不自动是给定多边形的单叶映射；边长、循环次序和简单性仍须核对。

\subsection{参数确定}
\label{b3:subsec:070304}

公式的常数 \(A\) 负责平移，\(C\) 负责旋转和缩放，预顶点的相对位置负责多边形的形状。上半平面自同构有三个实自由度，可固定三个预顶点；其余 \(n-3\) 个实参数通常须由边长比确定。边长条件写为
\[
 |w_{j+1}-w_j|=|C|\left|\int_{x_j}^{x_{j+1}}\prod_\ell(t-x_\ell)^{\alpha_\ell-1}\,\mathrm{d}t\right|.
\]
端点奇性可积，因为 \(\alpha_j-1>-1\)。相邻区间上的常辐角自动给出转角，积分模给边长。

\ProseParagraphBreak
预顶点可以很拥挤。长而细的多边形在圆盘或半平面参数中可能对应极靠近的预顶点，这称为数值计算中的拥挤现象。它不违反存在唯一性，但会使直接解边长方程不稳定。检查公式时，尺度自由度应先消除；否则把整体放大产生的等价参数当成多个不同解，会误判非唯一性。

\ProseParagraphBreak
参数数目也可以按边界资料理解。简单 \(n\) 边形的边方向由角度和整体旋转决定，边长还满足两个实闭合条件。缩放再消去一个自由度，留下 \(n-3\) 个独立的形状参数，恰与固定三个预顶点后的数目一致。这里只是维数计数，不能据此证明非线性边长方程在每一点有唯一解；对特殊对称形状、退化形状和数值分支，仍要检查映射对应的几何条件。

\ProseParagraphBreak
对于正三角形，可利用三个相等内角 \(\alpha_j=1/3\) 和半平面自同构把预顶点固定到三个选定实点。公式导数各有 \(-2/3\) 次端点奇性，边积分仍收敛，因为 \(\int_0^\varepsilon t^{-2/3}\,\mathrm{d}t=3\varepsilon^{1/3}\)。对于矩形，各内角为 \(1/2\)，四个预顶点在固定三个之后还剩一个实参数；这就是后一节椭圆积分中的 \(k\)。三角形无需额外形状参数，四边形需要一个，这是公式自由度的具体表现。

% !TeX root = ../../Book3.tex
\section{矩形与椭圆积分}
\label{b3:sec:0704}

矩形的四个角都相同，形状却仍有一个连续变化的参数。上一节的公式把这个自由度放在第四个预顶点的位置中；要找回矩形的宽高比，就必须实际计算相邻预顶点之间的积分。本节选择一个便于换元的归一化，先从积分构造出矩形映射，再研究边长随参数的变化。平方根在不同实区间的符号、无穷远对应的顶点和内部单叶性，都将在计算中逐一确定。

\subsection{矩形的共形映射}
\label{b3:subsec:070401}

将四个预顶点按正向次序放在
\[
 0,\quad 1,\quad a=\frac1{k^2},\quad\infty,
 \quad 0<k<1.
\]
这里用上半平面自同构固定了三个预顶点；写成 \(a=k^{-2}>1\)，是为了使后面的换元出现 \(1-k^2s^2\)。每个直角要求导数有 \(-1/2\) 次幂，因此三个有限顶点给出候选导数
\[
 F_k'(z)=\frac12q_k(z),\quad
 q_k(z)^2=\frac1{z(1-z)(1-k^2z)},\quad z\in\mathbb H.
\]
系数 \(1/2\) 只负责整体尺度，选它是因为 \(t=s^2\) 的微分恰有因子2。无穷远现在也是一个角点，不再是上一节公式中选定的边内点；候选导数在无穷远为 \(z^{-3/2}\) 阶，正好与这一变化相符。

\ProseParagraphBreak

平方根不能在每段实轴上独立取正值。对 \(z\in\mathbb H\)，分别使用满足
\[
 0<\arg z<\pi,\quad
 -\pi<\arg(1-z)<0,\quad
 -\pi<\arg(1-k^2z)<0
\]
的全纯对数，并定义
\[
 q_k(z)\coloneqq
 \exp\left[-\frac12\bigl(
 \operatorname{Log}z+\operatorname{Log}(1-z)
 +\operatorname{Log}(1-k^2z)\bigr)\right].
\]
三个因子在上半平面都不为零，各自的像又避开所选对数的割线，所以这给出一个全纯无零分支。在 \((0,1)\) 的上侧边界，三个辐角都趋于0，故 \(q_k\) 取正实值；这就确定了所需的第一条边向右的归一化。

\ProseParagraphBreak

上半平面单连通，\cref{b3:thm:02-cauchy-simply-connected}保证 \(q_k/2\) 有原函数。以边界点0为基点的写法
\[
 F_k(z)=\frac12\int_0^z
 \frac{\mathrm{d}t}{\sqrt{t(1-t)(1-k^2t)}}
\]
表示把这个原函数归一化为 \(F_k(0)=0\)，而不是未经说明地沿边界积分。为核对这个归一化存在，在0附近写
\[
 q_k(z)=z^{-1/2}A_k(z),\quad
 A_k(z)=\sum_{n=0}^{\infty}a_nz^n,\quad a_0=1.
\]
其中 \(A_k\) 在一个小圆盘全纯，平方根取上半平面分支。逐项积分得到局部原函数
\[
 \sum_{n=0}^{\infty}\frac{a_n}{2n+1}z^{n+1/2}
 =\sqrt z\bigl(1+O(z)\bigr).
\]
它从内部趋向0；任一全局原函数与它在小半圆盘中相差常数，减去该常数便得到唯一的 \(F_k\)。因此从0出发的可积端点和内部路径无关性都已得到说明。

\ProseParagraphBreak

现在观察实轴上侧的边界导数。按从0出发、经无穷远再回到0的循环次序，三个对数辐角的和及导数方向为
\[
\begin{array}{c|cccc}
 t&(0,1)&(1,a)&(a,\infty)&(-\infty,0)\\ \hline
 \arg t+\arg(1-t)+\arg(1-k^2t)
   &0&-\pi&-2\pi&\pi\\
 q_k(t)/|q_k(t)|&1&\mathrm{i}&-1&-\mathrm{i}
\end{array}
\]
例如越过1后，\(1-t\) 从下侧趋向负实轴，其辐角为 \(-\pi\)，倒平方根便给向上的方向。越过 \(a\) 后再出现一个 \(-\pi\)，两次相加为 \(-2\pi\)，对应向左。若先把这次辐角和改成0再取平方根，就会把第三条边误判为向右。负实轴上只有 \(t\) 的辐角为 \(\pi\)，故从 \(-\infty\) 走向0时导数向下。

\ProseParagraphBreak

在有限预顶点 \(x\in\{0,1,a\}\) 附近，其余因子都不为零，所以可写
\[
 q_k(z)=(z-x)^{-1/2}B_x(z),\quad B_x(x)\ne0,
\]
其中 \(B_x\) 在小圆盘全纯。与0处一样积分其幂级数，得到
\[
 F_k(z)=c_x+(z-x)^{1/2}C_x(z),\quad C_x(x)=B_x(x)\ne0.
\]
于是定义 \(F_k(x)=c_x\) 后，\(F_k\) 从上半平面的任意方向连续到 \(x\)，差值为 \(O(|z-x|^{1/2})\)。平方根把半圆盘的角度 \(\pi\) 变为 \(\pi/2\)，说明这三个顶点确为直角；导数发散不妨碍函数有有限边界值。

\ProseParagraphBreak

无穷远还需要单独处理，否则正、负实轴得到的两段边界可能尚未接上。当 \(|z|\ge2/k^2\) 时，
\[
 |1-z|\ge\frac{|z|}{2},\quad
 |1-k^2z|\ge\frac{k^2|z|}{2},\quad
 |F_k'(z)|\le\frac1{k|z|^{3/2}}.
\]
正、负实轴上的尾积分因此都收敛。连接 \(R\) 与 \(-R\) 的上半圆弧长为 \(\pi R\)，其上的积分模至多为 \(\pi/(k\sqrt R)\)，故两侧极限相同，记为 \(L_k\)。对任意 \(z\in\mathbb H\)、\(|z|=R\)，先沿同一半圆从 \(R\) 走到 \(z\)，再估计实轴尾积分，得到
\[
 |F_k(z)-L_k|\le\frac{\pi+2}{k\sqrt{|z|}}.
\]
这个估计证明从内部任何方向趋于无穷远都得到 \(L_k\)，并非只核对两个实轴端点。

\ProseParagraphBreak

第四个角也可以在有限坐标中看清。令 \(\zeta=-1/z\)，它仍把上半平面映到上半平面。由上面的分支在大正实数处为负实值，可将无穷远展开写成
\[
 q_k(z)=-\frac1kz^{-3/2}\bigl(1+O(1/z)\bigr).
\]
积分并用已确定的 \(L_k\) 归一化，得到
\[
 F_k(-1/\zeta)
 =L_k-\frac{\mathrm{i}}k\sqrt\zeta+O(|\zeta|^{3/2}).
\]
这里 \(0<\arg\zeta<\pi\)。非零平方根首项再次给出直角。到此，四段边界的方向与连续接合已经确定；它们的长度和所围区域仍须由积分计算。

\subsection{椭圆积分}
\label{b3:subsec:070402}

第一条边的积分在 \(t=s^2\) 下成为一个经常出现的标准量。对 \(0\le k<1\)，定义\term[diyilei-wanquan-tuoyuan-jifen]{第一类完全椭圆积分}[complete elliptic integral of the first kind]
\[
 K(k)\coloneqq\int_0^1
 \frac{\mathrm{d}s}{\sqrt{(1-s^2)(1-k^2s^2)}}
 =\int_0^{\pi/2}\frac{\mathrm{d}\theta}
 {\sqrt{1-k^2\sin^2\theta}}.
\]
\symidx{\(K(k)\)：第一类完全椭圆积分}
第二个表达式来自 \(s=\sin\theta\)。对固定 \(k<1\)，被积函数至多为
\((1-k^2)^{-1/2}(1-s^2)^{-1/2}\)，所以积分有限且为正。还记
\[
 k'\coloneqq\sqrt{1-k^2}.
\]
撇号在这里表示与 \(k\) 配对的参数，不是求导；再次进行同一变换会回到 \(k\)。

\Needspace{6\baselineskip}
\begin{proposition}[椭圆积分与矩形]\label{b3:prop:07-elliptic-rectangle}
设 \(0<k<1\)、\(k'=\sqrt{1-k^2}\)，并令
\[
 F_k(z)=\frac12\int_0^z
 \frac{\mathrm{d}t}{\sqrt{t(1-t)(1-k^2t)}},\quad z\in\mathbb H,
\]
其中倒平方根在上半平面全纯，在 \((0,1)\) 的上侧取正值，原函数归一化为 \(F_k(0)=0\)。则 \(F_k\) 双全纯映到矩形
\[
 R_k\coloneqq\{w\in\mathbb C:
 0<\operatorname{Re}w<K(k),\quad
 0<\operatorname{Im}w<K(k')\}.
\]
它连续延到扩充实轴，依次将 \(0,1,k^{-2},\infty\) 映为
\[
 0,\quad K(k),\quad K(k)+\mathrm{i}K(k'),\quad
 \mathrm{i}K(k'),
\]
并给出扩充实轴与矩形边界之间的同胚。
\end{proposition}
\begin{proof}
先计算四条边长。对 \((0,1)\)，导数为正实数；取 \(t=s^2\)、\(\mathrm{d}t=2s\,\mathrm{d}s\)，得到
\[
 \ell_1=\frac12\int_0^1
 \frac{\mathrm{d}t}{\sqrt{t(1-t)(1-k^2t)}}
 =\int_0^1\frac{\mathrm{d}s}
 {\sqrt{(1-s^2)(1-k^2s^2)}}=K(k).
\]
所以 \(F_k(1)=K(k)\)。

第二条边向上。先仍取 \(t=s^2\)，则其正长度为
\[
 \ell_2=\frac12\int_1^{1/k^2}
 \frac{\mathrm{d}t}{\sqrt{t(t-1)(1-k^2t)}}
 =\int_1^{1/k}\frac{\mathrm{d}s}
 {\sqrt{(s^2-1)(1-k^2s^2)}}.
\]
为把两个端点变成0和1，令
\[
 s=(1-k'^2u^2)^{-1/2},\quad 0\le u\le1.
\]
这里 \(s=1\) 对应 \(u=0\)，\(s=1/k\) 对应 \(u=1\)。逐项计算有
\[
 \begin{aligned}
 s^2-1&=\frac{k'^2u^2}{1-k'^2u^2},&
 1-k^2s^2&=\frac{k'^2(1-u^2)}{1-k'^2u^2},\\
 \mathrm{d}s&=\frac{k'^2u}{(1-k'^2u^2)^{3/2}}\,\mathrm{d}u.
 \end{aligned}
\]
在 \(0<u<1\) 上取正平方根并约去 \(k'^2u\)，得到
\[
 \frac{\mathrm{d}s}{\sqrt{(s^2-1)(1-k^2s^2)}}
 =\frac{\mathrm{d}u}{\sqrt{(1-u^2)(1-k'^2u^2)}}.
\]
两个端点的可积性由右式一并给出，因此 \(\ell_2=K(k')\)，且
\(F_k(k^{-2})=K(k)+\mathrm{i}K(k')\)。

第三条边向左，长度是
\[
 \ell_3=\frac12\int_{1/k^2}^{\infty}
 \frac{\mathrm{d}t}{\sqrt{t(t-1)(k^2t-1)}}.
\]
令 \(t=1/(k^2s^2)\)，则 \(t\) 从 \(k^{-2}\) 增至无穷时，\(s\) 从1减至0，而且
\[
 \begin{aligned}
 \mathrm{d}t&=-\frac{2}{k^2s^3}\,\mathrm{d}s,\\
 t-1&=\frac{1-k^2s^2}{k^2s^2},\quad
 k^2t-1=\frac{1-s^2}{s^2},\\
 \sqrt{t(t-1)(k^2t-1)}
 &=\frac{\sqrt{(1-s^2)(1-k^2s^2)}}{k^2s^3}.
 \end{aligned}
\]
故 \(\ell_3=K(k)\)。这一边从 \(K(k)+\mathrm{i}K(k')\) 向左走过 \(K(k)\)，到达前一小节确定的无穷远极限
\(L_k=\mathrm{i}K(k')\)。

最后，从负无穷走向0的第四条边向下，其长度为
\[
 \ell_4=\frac12\int_{-\infty}^0
 \frac{\mathrm{d}t}{\sqrt{(-t)(1-t)(1-k^2t)}}.
\]
取 \(t=-s^2/(1-s^2)\)，原积分的两端依次对应 \(s=1,0\)。有
\[
 \begin{aligned}
 -t&=\frac{s^2}{1-s^2},\quad
 1-t=\frac1{1-s^2},\quad
 1-k^2t=\frac{1-k'^2s^2}{1-s^2},\\
 \mathrm{d}t&=-\frac{2s}{(1-s^2)^2}\,\mathrm{d}s,\\
 \sqrt{(-t)(1-t)(1-k^2t)}
 &=\frac{s\sqrt{1-k'^2s^2}}{(1-s^2)^{3/2}}.
 \end{aligned}
\]
约分并反转积分限，得到
\[
 \ell_4=\int_0^1\frac{\mathrm{d}s}
 {\sqrt{(1-s^2)(1-k'^2s^2)}}=K(k').
\]
所以第四边从 \(\mathrm{i}K(k')\) 回到0。四个实区间上导数的模始终为正，各边被严格单调地走过一次；结合端点连续性，扩充实轴恰好按正向描出这个简单矩形边界。

剩下的是内部单叶性。不能直接在实轴及无穷远组成的边界上套用要求全纯邻域的辐角原理。改用 Cayley 坐标
\[
 T(\zeta)=\mathrm{i}\frac{1+\zeta}{1-\zeta},\quad
 g(\zeta)=F_k(T(\zeta)),\quad \zeta\in\mathbb D.
\]
有限角点及无穷远的估计说明 \(g\) 连续到闭圆盘，其中 \(g(1)=L_k\)。当 \(\zeta=e^{\mathrm{i}t}\)、\(0<t<2\pi\) 时，
\(T(\zeta)=-\cot(t/2)\) 从负无穷严格增加到正无穷。因此
\(\gamma_1(t)=g(e^{\mathrm{i}t})\)
按正向描出矩形边界，从左上角出发先向下，再向右、向上、向左。

固定 \(w\notin\partial R_k\)，令 \(d=\operatorname{dist}(w,\partial R_k)>0\)。闭圆盘上的一致连续性给
\[
 \sup_{0\le t\le2\pi}
 |g(re^{\mathrm{i}t})-g(e^{\mathrm{i}t})|\longrightarrow0
 \quad(r\uparrow1).
\]
于是存在 \(r_0<1\)，使对每个 \(r_0<r<1\)，路径
\(\gamma_r(t)=g(re^{\mathrm{i}t})\)
与 \(\gamma_1\) 的距离处处小于 \(d/2\)。两条路径之间的直线同伦始终避开 \(w\)，所以连续辐角的总增量不变，绕数相同。\Cref{b3:thm:03-argument-principle}现在可用于 \(|\zeta|<r\)：函数 \(g-w\) 在该闭圆盘的邻域全纯，且边界无零点、内部无极点。因此其内部零点总重数为
\[
 N_r(w)=\operatorname{Ind}(\gamma_r,w)
 =\operatorname{Ind}(\gamma_1,w)
 =\begin{cases}
 1,&w\in R_k,\\
 0,&w\notin\overline{R_k}.
 \end{cases}
\]
这里矩形的正向边界对内部点绕一周，对外部点绕零周。

若 \(w\in R_k\)，先取一个 \(r>r_0\)，便得到至少一个原像；若全圆盘中另有原像，或一个零点的重数大于1，取更大的 \(r<1\) 将它们包括在内，就与 \(N_r(w)=1\) 矛盾。若 \(w\notin\overline{R_k}\)，任何可能的原像也会包含在某个这样的紧盘内，与零计数矛盾。最后，内部点不能映到 \(\partial R_k\)：否则开映射定理给出的像邻域必含矩形外部点，违反刚得到的结论。故 \(F_k\) 的像恰为 \(R_k\)，并且一一对应。又因 \(F_k'=q_k/2\ne0\)，局部全纯逆拼成全局全纯逆映射，证明完成。
\end{proof}

四边长度的换元具有不同任务。第一边给出 \(K(k)\) 的定义，第二边引出 \(k'\)，第三、第四边独立核对闭合条件；最后的绕数则把边界上一圈转化为内部一个原像。导数处处非零只保证局部可逆，不能单独代替这一步全局计数。

\subsection{矩形模量}
\label{b3:subsec:070403}

对宽为 \(A>0\)、高为 \(B>0\) 的矩形，指定左、右两条边为需要连接的对边时，本节取\term[juxing-moliang]{矩形模量}[modulus of a rectangle]
\[
 m\coloneqq\frac AB.
\]
这是带有方向约定的形状参数。在上面的映射中，底边宽度为 \(K(k)\)，左边高度为 \(K(k')\)，所以
\[
 m(k)=\frac{K(k)}{K(k')}.
\]
若改为连接上、下两边，宽高角色交换，所得参数为 \(1/m\)。下一章会用极值长度说明这个比值为何成为一般标记四边形的共形不变量；本节先证明每个正比值都能由积分唯一产生。

\Needspace{6\baselineskip}
\begin{proposition}[矩形宽高比的参数化]\label{b3:prop:07-elliptic-aspect}
对 \(0\le t<1\)，令
\[
 K(t)=\int_0^1\frac{\mathrm{d}s}
 {\sqrt{(1-s^2)(1-t^2s^2)}}.
\]
则 \(K\) 在 \([0,1)\) 连续且严格递增，\(K(0)=\pi/2\)，并在 \(t\uparrow1\) 时趋于无穷。对 \(0<k<1\) 置 \(k'=\sqrt{1-k^2}\)，则
\(m(k)=K(k)/K(k')\)
是从 \((0,1)\) 到 \((0,\infty)\) 的连续严格递增双射，且
\[
 m(k')=\frac1{m(k)}.
\]
因此每个指定左、右边的矩形宽高比唯一确定一个 \(k\)。
\end{proposition}
\begin{proof}
若 \(0\le t_1<t_2<1\)，则对每个 \(0<s<1\)，第二个根号因子在 \(t_2\) 时严格较小，被积函数严格较大。两者均可积，故 \(K(t_2)>K(t_1)\)。对任意 \(b<1\)，当 \(0\le t\le b\) 时有共同可积控制
\[
 \frac1{\sqrt{(1-s^2)(1-t^2s^2)}}
 \le\frac1{\sqrt{1-b^2}\sqrt{1-s^2}}.
\]
所以控制收敛给出 \([0,b]\) 上的连续性，\(b\) 任意便得到 \([0,1)\) 上连续。取 \(t=0\) 后作 \(s=\sin\theta\)，有 \(K(0)=\pi/2\)。

当 \(t\uparrow1\)，被积函数递增趋于 \(1/(1-s^2)\)，而后者在 \(s=1\) 不可积。这个无穷极限也可直接由截断核对：给定 \(M>0\)，先取 \(c<1\)，使
\(\int_0^c(1-s^2)^{-1}\,\mathrm{d}s>2M\)；在固定区间 \([0,c]\) 上，被积函数一致趋于该极限，故 \(t\) 充分接近1时，\(K(t)\ge\int_0^c((1-s^2)(1-t^2s^2))^{-1/2}\,\mathrm{d}s>M\)。

随 \(k\) 严格增大，\(K(k)\) 严格增大，\(k'\) 严格减小，\(K(k')\) 严格减小。二者为正，故其比严格递增，并由连续性得到连续。两端分别满足
\[
 \lim_{k\downarrow0}m(k)=0,\quad
 \lim_{k\uparrow1}m(k)=\infty.
\]
介值定理给满射，严格递增给单射。最后，\(\sqrt{1-(k')^2}=k\)，交换分子与分母即得互反关系。
\end{proof}

例如 \(k=1/\sqrt2\) 时有 \(k=k'\)，矩形为正方形；反之，若像为正方形，则 \(K(k)=K(k')\)，严格单调性迫使 \(k=k'\)，所以只有这一参数。给定任意宽 \(A\)、高 \(B\)，先解唯一的方程 \(m(k)=A/B\)，再把 \(F_k\) 乘以
\[
 c=\frac{A}{K(k)}=\frac{B}{K(k')},
\]
就得到以0为左下角、宽 \(A\)、高 \(B\) 的矩形映射。形状参数与整体尺度在这两个步骤中分别确定。

\ProseParagraphBreak

还可定量解释细长矩形的预顶点为什么难以分辨。直接估计给
\[
 \frac\pi2\le K(k)\le\frac\pi{2\sqrt{1-k^2}},
\]
但右端在 \(k\uparrow1\) 时是幂次发散，远大于实际增长。为了得到正确的对数尺度，换用 \(u=\tan\theta\)，记 \(\varepsilon=k'\)。因为
\[
 \mathrm{d}\theta=\frac{\mathrm{d}u}{1+u^2},\quad
 1-k^2\sin^2\theta=\frac{1+\varepsilon^2u^2}{1+u^2},
\]
得到
\[
 K(k)=\int_0^\infty
 \frac{\mathrm{d}u}{\sqrt{(1+u^2)(1+\varepsilon^2u^2)}}.
\]
在 \(1\le u\le1/\varepsilon\) 上，两个平方根分别接近 \(u\) 与1；因此增长来自近似 \(1/u\) 的一段积分，而两端只贡献有界的量。

\Needspace{6\baselineskip}
\begin{proposition}[完全椭圆积分的对数界]\label{b3:prop:07-elliptic-log-bounds}
设 \(0<k<1\)、\(\varepsilon=\sqrt{1-k^2}\)，并令
\(K(k)=\int_0^1((1-s^2)(1-k^2s^2))^{-1/2}\,\mathrm{d}s\)。则
\[
 \log\frac1\varepsilon-\frac12
 \le K(k)\le\log\frac1\varepsilon+2.
\]
特别地，\(K(k)=\log(1/k')+O(1)\) 当 \(k\uparrow1\)，其中误差有与 \(k\) 无关的绝对界。
\end{proposition}
\begin{proof}
将刚得到的 \(u\) 积分分为 \([0,1]\)、\([1,1/\varepsilon]\) 与 \([1/\varepsilon,\infty)\)。第一段的被积函数至多为1，最后一段至多为 \(1/(\varepsilon u^2)\)，因此这两段各不超过1。中间一段可写为
\[
 \int_1^{1/\varepsilon}\frac1u
 \frac{\mathrm{d}u}{\sqrt{(1+u^{-2})(1+\varepsilon^2u^2)}}
 \le\int_1^{1/\varepsilon}\frac{\mathrm{d}u}{u}
 =\log\frac1\varepsilon.
\]
相加即得上界。

为估计中间一段少于对数的部分，先注意对 \(v\ge0\)，
\[
 0\le1-(1+v)^{-1/2}
 =\frac{v}{\sqrt{1+v}\,(1+\sqrt{1+v})}
 \le\frac v2.
\]
若 \(a,b\ge0\)，将两个因子的差分开，便有
\[
 0\le1-\frac1{\sqrt{(1+a)(1+b)}}
 \le\frac{a+b}{2}.
\]
代入 \(a=u^{-2}\)、\(b=\varepsilon^2u^2\)，中间一段与 \(\log(1/\varepsilon)\) 之差至多为
\[
 \frac12\int_1^{1/\varepsilon}
 (u^{-3}+\varepsilon^2u)\,\mathrm{d}u
 =\frac{1-\varepsilon^2}{2}\le\frac12.
\]
其余两段非负，所以舍去它们仍得到所述下界。
\end{proof}

这一估计已经足以证明指数级的拥挤，而无需先求更精细的渐近常数。令 \(k\ge1/\sqrt2\)，则 \(0<\varepsilon\le1/\sqrt2\)、\(m=m(k)\ge1\)。由最初的上下界和
\[
 \frac1{\sqrt{1-\varepsilon^2}}-1
 =\frac{\varepsilon^2}
 {\sqrt{1-\varepsilon^2}\,(1+\sqrt{1-\varepsilon^2})}
 \le\varepsilon^2
\]
得到
\[
 0\le K(\varepsilon)-\frac\pi2
 \le\frac\pi2\varepsilon^2.
\]
记 \(L=\log(1/\varepsilon)\)。由于
\[
 L-\frac\pi2m
 =L-K(k)+m\left(K(\varepsilon)-\frac\pi2\right),
\]
第一项在 \([-2,1/2]\) 内，而第二项满足
\[
 0\le m\left(K(\varepsilon)-\frac\pi2\right)
 \le\varepsilon^2(L+2)
 \le\frac1{2e}+1.
\]
最后一步用了 \(\varepsilon^2\log(1/\varepsilon)\le1/(2e)\) 及 \(2\varepsilon^2\le1\)。因此有具体的统一误差界
\[
 \left|\log\frac1\varepsilon-\frac\pi2m\right|\le2.
\]
这里保留了 \(K(\varepsilon)-\pi/2\) 的定量控制；仅知道它趋于0，不足以在乘上不断增大的 \(m\) 后仍声称误差有界。

\ProseParagraphBreak

在固定预顶点 \(0,1,\infty\) 的归一化中，另外两个有限预顶点 \(1\) 与 \(k^{-2}\) 的间隙为
\[
 \Delta=\frac1{k^2}-1=\frac{\varepsilon^2}{1-\varepsilon^2}.
\]
当 \(k\ge1/\sqrt2\) 时，\(\varepsilon^2\le\Delta\le2\varepsilon^2\)。将前式指数化便得到
\[
 e^{-4}e^{-\pi m}\le\Delta\le2e^4e^{-\pi m}.
\]
常数并不精细，但结论很明确：宽高比增大时，归一化后的两个预顶点以指数尺度靠近。因而即使矩形的边长输入仍容易辨认，求解预顶点位置也可能需要很多有效数字。对很高而窄的矩形，交换所指定的边，同时将 \(k\) 换成 \(k'\)，就把同一估计用于 \(1/m\)。这项比较针对已经固定三个点的坐标，不把可任意缩放的实轴距离当成共形不变量。

\subsection{特殊函数联系}
\label{b3:subsec:070404}

完全积分记录一条完整边的长度；若终点也作为变量，就得到相应的不完全积分。在 \(s=0\) 的邻域取平方根值为1，定义
\[
 J_k(s)\coloneqq\int_0^s
 \frac{\mathrm{d}u}{\sqrt{(1-u^2)(1-k^2u^2)}}.
\]
例如在 \(|s|<1\) 内，两个因子均无零点，可取从0开始的全纯平方根，积分与路径无关。由于 \(J_k(0)=0\)、\(J_k'(0)=1\)，\cref{b3:thm:01-local-inverse}给出0附近唯一的全纯逆。它记作 \(\operatorname{sn}(w,k)\)，是\term[jacobi-tuoyuan-hanshu]{Jacobi 椭圆函数}[Jacobi elliptic function] \(\operatorname{sn}\) 在0附近的定义，并满足
\[
 J_k(\operatorname{sn}(w,k))=w,\quad
 \operatorname{sn}(0,k)=0,\quad
 \frac{\partial\operatorname{sn}}{\partial w}(0,k)=1.
\]
\symidx{\(\operatorname{sn}(w,k)\)：不完全椭圆积分在0附近的归一化逆函数}
此时的定义域只是0的一个邻域，尚未赋予它全平面上的值或周期。

\ProseParagraphBreak

链式法则给这一个局部逆的微分关系。若暂写 \(S(w)=\operatorname{sn}(w,k)\)，则
\[
 S'(w)=\sqrt{(1-S(w)^2)(1-k^2S(w)^2)},\quad
 (S'(w))^2=(1-S(w)^2)(1-k^2S(w)^2).
\]
第一式的平方根在0取值1。第二式虽然省去了分支记号，却不能代替这个初值选择：仅有平方后的等式会丢掉导数的符号。又因 \(J_k'\) 为偶函数且 \(J_k(0)=0\)，\(J_k\) 为奇函数，逆函数的唯一性进一步给 \(S(-w)=-S(w)\)。

\ProseParagraphBreak

局部逆与矩形映射之间还隔着一次平方。对充分小且位于第一象限的 \(s\)，有 \(s^2\in\mathbb H\)，换元 \(t=u^2\) 并使用同一分支可得
\[
 F_k(s^2)=J_k(s).
\]
因此在矩形左下角附近，
\[
 F_k^{-1}(w)=\operatorname{sn}(w,k)^2.
\]
这里不能把 \(F_k^{-1}\) 直接写成 \(\operatorname{sn}\)。实际上，\(F_k^{-1}\) 在整个 \(R_k\) 内取值于上半平面，故其取值于第一象限的平方根
\(\sqrt{F_k^{-1}(w)}\)
把这个局部 \(\operatorname{sn}\) 继续到了矩形内部。接下来可以利用直边反射研究相邻矩形中的延拓，但反复反射时还要跟踪平方根的符号、顶点附近的延拓和可能出现的极点。继续延拓并比较相邻矩形的坐标，就会引出周期问题；后面的椭圆函数理论将研究这种局部逆的全平面亚纯表示及双周期性。

\ProseParagraphBreak

参数为0的退化也先从局部积分理解。令 \(k=0\)，则
\[
 J_0(s)=\int_0^s\frac{\mathrm{d}u}{\sqrt{1-u^2}}
 =\arcsin s
\]
是0附近取值为0的反正弦分支，因而 \(\operatorname{sn}(w,0)=\sin w\) 在0附近成立。对上半平面则可直接构造
\[
 F_0(z)=\frac12\int_0^z\frac{\mathrm{d}t}{\sqrt{t(1-t)}}
 =\arcsin\sqrt z,
\]
其中右端先在0附近的上半邻域确定，再由左端的原函数延到整个 \(\mathbb H\)。这避免了在远离0时任意选取反正弦的其他分支。

\ProseParagraphBreak

还应说明 \(F_0\) 确实是 \(F_k\) 的极限。对每个 \(E\Subset\mathbb H\)，因子 \((1-k^2z)^{-1/2}\) 在 \(E\) 上一致趋于1，所以 \(F_k'\to F_0'\) 局部一致。在固定路径 \(t=\mathrm{i}u\)、\(0<u\le1\) 上，积分核的模至多为 \(u^{-1/2}\)，它可积，故控制收敛给 \(F_k(\mathrm{i})\to F_0(\mathrm{i})\)。再沿从 \(\mathrm{i}\) 到 \(z\in E\) 的直线段积分：这些线段合在一个仍位于上半平面的紧集中，长度有共同上界，因此导数的局部一致收敛给出
\[
 F_k\longrightarrow F_0
 \quad\text{在 }\mathbb H\text{ 的每个紧集上一致}.
\]
这里没有把 \(k=0\) 代入四个仍互异的预顶点资料；实际上 \(k^{-2}\) 已经逃向无穷，原来的两个顶点正在合并。

\ProseParagraphBreak

因为 \(K(k)\to\pi/2\)、\(K(k')\to\infty\)，并且 \(F_k(\mathbb H)=R_k\)，极限的像落在
\(\{0\le\operatorname{Re}w\le\pi/2,\ \operatorname{Im}w\ge0\}\)
中。\(F_0'\ne0\)，开映射定理排除等号边界，所以其像包含于半无限条带
\[
 S_0\coloneqq\{w\in\mathbb C:
 0<\operatorname{Re}w<\pi/2,\quad\operatorname{Im}w>0\}.
\]
要确认没有遗漏条带中的点，可以直接使用逆式。在0附近有
\(\sin^2F_0(z)=z\)，恒等定理将它推广到整个上半平面。而对 \(w=u+\mathrm{i}v\in S_0\)，
\[
 \operatorname{Im}(\sin^2w)
 =\frac12\sin(2u)\sinh(2v)>0,
\]
故 \(\sin^2\) 把 \(S_0\) 映入 \(\mathbb H\)。若 \(w_1,w_2\in S_0\) 且 \(\sin^2w_1=\sin^2w_2\)，恒等式
\[
 \sin^2w_1-\sin^2w_2
 =\sin(w_1+w_2)\sin(w_1-w_2)
\]
说明 \(w_1+w_2\) 或 \(w_1-w_2\) 为 \(\pi\) 的整数倍。前者因虚部为正不可能；后者的实部在 \((-\pi/2,\pi/2)\)，只能为0。因此 \(\sin^2\) 在 \(S_0\) 单射。给定 \(w\in S_0\)，令 \(z=\sin^2w\)，则 \(F_0(z)\in S_0\) 且与 \(w\) 具有相同的平方正弦，由单射性得 \(F_0(z)=w\)。所以 \(F_0\) 恰好双全纯映到 \(S_0\)，逆映射为 \(\sin^2w\)。

\ProseParagraphBreak

这个极限保留了底边的两个顶点，而把上边推向无穷，因而得到的是半无限条带。对每个固定的 \(k>0\)，像仍为四边有限的矩形。局部积分的收敛、特殊函数的退化和边界预顶点的合并在这里相互吻合；全局周期如何从非退化矩形产生，则需要\chapref{b3:ch:22}建立的格与亚纯函数理论。


\begin{chaptersummary}
普通边界上的连续性、边界单射性与导数存在性是不同结论。素端提供任意单连通区域的圆周式边界，Jordan 区域使它与平面边界一致。多边形角度控制导数幂次，边长控制预顶点参数；矩形的唯一形状参数由椭圆积分比给出。
\end{chaptersummary}

\BookExercises

% 本章题组自拟；以正文已建立的工具设计计算、反例、构造与综合训练。
\begin{exercise}\label{b3:ex:07-1}
说明圆盘裂缝 \(\mathbb D\setminus[0,1)\) 中，裂缝内点从两侧趋近为何必须对应不同素端。
\end{exercise}
\begin{solution}
在该点附近取不跨裂缝的上下两个小半圆横割，并沿各自一侧取嵌套分支。充分晚分支分别位于裂缝上方与下方，互不包含，因此链不互相从属。印象都为选定点，所以相同印象不推出素端相同。
\end{solution}

\begin{exercise}\label{b3:ex:07-2}
若 Jordan 区域映射连续延到边界，是否必有连续边界导数？用多边形角点说明。
\end{exercise}
\begin{solution}
不必。内角 \(\alpha\pi\) 对应 \(f'(z)\sim C(z-x)^{\alpha-1}\)。凸角 \(\alpha<1\) 时导数模趋于无穷；凹角 \(\alpha>1\) 时导数趋于0。边界同胚不要求导数有非零有限极限。
\end{solution}

\begin{exercise}\label{b3:ex:07-3}
设 \(g\) 在上半圆盘全纯且连续到直径，直径上满足 \(\operatorname{Im}g=c\)。写出其反射延拓。
\end{exercise}
\begin{solution}
先令 \(h=g-\mathrm{i}c\)，直径上实值。下半盘取 \(h(z)=\overline{h(\overline z)}\)，所以 \(g(z)=\overline{g(\overline z)}+2\mathrm{i}c\)。Morera 拼接证明保证全纯，不是仅凭连续性断言。
\end{solution}

\begin{exercise}\label{b3:ex:07-4}
求三角形三个有限预顶点形式中导数的指数和，并解释无穷远的行为。
\end{exercise}
\begin{solution}
内角和为 \(\pi\)，所以 \(\sum\alpha_j=1\)，指数和 \(\sum(\alpha_j-1)=-2\)。于是 \(f'(z)\sim Cz^{-2}\)，积分在无穷远有有限极限且以 \(1/z\) 为局部一阶坐标，符合无穷远对应边内点。
\end{solution}

\begin{exercise}\label{b3:ex:07-5}
说明 Schwarz--Christoffel 边长积分在有限预顶点处为何可积，凹角是否例外。
\end{exercise}
\begin{solution}
奇性阶为 \(\alpha_j-1>-1\)，因此 \(\int_0^\varepsilon t^{\alpha_j-1}\,\mathrm{d}t=\varepsilon^{\alpha_j}/\alpha_j<\infty\)。凹角的指数反而为正，导数趋于0，更不存在积分发散。
\end{solution}

\begin{exercise}\label{b3:ex:07-6}
验证 \(K(k')\) 的换元公式，并求正方形所对应的 \(k\)。
\end{exercise}
\begin{solution}
在边长积分取 \(s=(1-k'^2u^2)^{-1/2}\)，端点1与 \(1/k\) 对应0与1，约分后正是 \(K(k')\)。正方形要求 \(K(k)=K(k')\)，严格单调性迫使 \(k=k'\)，即 \(k=1/\sqrt2\)。
\end{solution}

\begin{exercise}\label{b3:ex:07-7}
为什么只指定四个角为直角不足以确定一个矩形的共形标记？
\end{exercise}
\begin{solution}
所有矩形的四角相同，却有不同边长比。固定三个预顶点后，第四个预顶点留下一个实参数；公式中的 \(k\) 决定 \(K(k)/K(k')\)。改变整体尺度不改变该比，无法消除这个参数。
\end{solution}

\begin{exercise}\label{b3:ex:07-8}
设 \(f,g\) 是同一 Jordan 区域的两个圆盘映射，且三个循环有序圆周点的边界像分别相同。证明 \(f=g\)。
\end{exercise}
\begin{solution}
复合 \(g^{-1}\circ f\) 是圆盘自同构，并通过边界同胚固定这三个圆周点。分式变换固定点方程是次数至多2的多项式；若有三个不同固定点，方程恒零，故自同构恒等。于是两映射相同。
\end{solution}
